\clearpage
\section{사차 갈루아군의 다섯 표준 예제}
\label{sec:quartic-worked-examples}

\begin{learninggoal}
삼차분해식·판별식·분해체 차수를 함께 사용하여 $V_4,C_4,D_4,A_4,S_4$의 예를 완전히 계산한다. 하나의 불변량만으로 결론을 서두르지 않는 계산 습관을 익힌다.
\end{learninggoal}

\subsection{클라인 네 군: 두 독립인 제곱근}

\begin{example}
\label{ex:quartic-v4}
\[
  f=X^4-10X^2+1
\]
의 갈루아군은 $V_4$이다.
\end{example}

\begin{proof}
$\alpha=\sqrt3+\sqrt2$라 두면
\[
  \alpha^{-1}=\sqrt3-\sqrt2
\]
이고
\[
  \sqrt3=\frac{\alpha+\alpha^{-1}}2,
  \qquad
  \sqrt2=\frac{\alpha-\alpha^{-1}}2.
\]
따라서
\[
  \Q(\alpha)=\Q(\sqrt2,\sqrt3)
\]
이고 이 체의 차수는 $4$이다. $f$의 근은
\[
  \pm(\sqrt3+\sqrt2),
  \qquad
  \pm(\sqrt3-\sqrt2)
\]
이므로 이 체가 분해체이다. 갈루아군은 두 제곱근의 부호를 독립적으로 바꾸는 군
\[
  C_2\times C_2\cong V_4
\]
이다.

계산 절차로도 확인할 수 있다. 여기서 $p=-10,q=0,r=1$이고
\[
  R_f(Z)=Z^3+20Z^2+96Z
  =Z(Z+8)(Z+12)
\]
는 완전히 분해된다. 또한
\[
  \Disc(f)=147456=384^2.
\]
\end{proof}

\subsection{순환 사차군: 이미 한 근이 모든 근을 담는 경우}

\begin{example}
\label{ex:quartic-c4}
\[
  f=X^4-4X^2+2
\]
의 갈루아군은 $C_4$이다.
\end{example}

\begin{proof}
$f$는 소수 $2$에 대한 아이젠슈타인 판정법으로 기약이다. $\alpha$를
\[
  \alpha^2=2+\sqrt2
\]
를 만족하는 근으로 택한다. 그러면
\[
  \sqrt2=\alpha^2-2\in\Q(\alpha).
\]
다른 양의 제곱근을 $\beta=\sqrt{2-\sqrt2}$라 하면
\[
  \alpha\beta=\sqrt2,
  \qquad
  \beta=\frac{\sqrt2}{\alpha}\in\Q(\alpha).
\]
따라서 네 근 $\pm\alpha,\pm\beta$가 모두 $\Q(\alpha)$ 안에 있고, 기약성에서
\[
  [\Q(\alpha):\Q]=4.
\]
즉 분해체 차수가 $4$이다.

삼차분해식은
\[
  R_f(Z)=Z(Z^2+8Z+8)
\]
로 $1+2$형이고, 판별식은
\[
  \Disc(f)=2048
\]
로 유리수의 제곱이 아니다. 따라서 위수 $4$의 갈루아군은 $V_4$가 아니라
\[
  C_4
\]
이다. 실제로 $R_f$의 분해체는 $F=\Q(\sqrt2)$이고 $[L:F]=2$이다.
\end{proof}

\subsection{이면체 사차군: 네제곱근과 허수단위}

\begin{example}
\label{ex:quartic-d4}
\[
  f=X^4-2
\]
의 갈루아군은 $D_4$이다.
\end{example}

\begin{proof}
$a=\sqrt[4]{2}>0$라 두면 근은
\[
  a,\quad ia,\quad -a,\quad -ia
\]
이고 분해체는
\[
  L=\Q(a,i)
\]
이다. $f$는 아이젠슈타인 판정법으로 기약이므로 $[\Q(a):\Q]=4$이다. $\Q(a)$는 실수체이므로 $i\notin\Q(a)$이고
\[
  [L:\Q]=8.
\]
네 근 위에서 추이적으로 작용하는 $S_4$의 위수 $8$ 부분군은 $D_4$이다.

한편 $p=q=0,r=-2$이므로
\[
  R_f(Z)=Z(Z^2+8).
\]
그 분해체는 $F=\Q(\sqrt{-2})$이고
\[
  [L:F]=4,
\]
여기서도 $D_4$ 판정이 나온다. 판별식은
\[
  \Disc(f)=-2048.
\]
\end{proof}

\subsection{교대 사차군: 기약 분해식과 제곱 판별식}

\begin{example}
\label{ex:quartic-a4}
\[
  f=X^4-7X^2-3X+1
\]
의 갈루아군은 $A_4$이다.
\end{example}

\begin{proof}
유리근은 없다. 만약 $f$가 두 일계수 이차식의 곱이라면 가우스 보조정리에 의해
\[
  f=(X^2+aX+b)(X^2-aX+d)
  \qquad(a,b,d\in\Z)
\]
로 쓸 수 있다. 상수항에서 $bd=1$이므로 $b=d=1$ 또는 $b=d=-1$이다. 그러면 일차항 $a(d-b)$가 $0$이 되어 실제 일차항 $-3X$와 모순이다. 따라서 $f$는 기약이다.

삼차분해식은
\[
  R_f(Z)=Z^3+14Z^2+45Z+9.
\]
유리근 정리에 따라 가능한 근 $\pm1,\pm3,\pm9$를 대입하면 어느 것도 근이 아니므로 이 삼차식은 기약이다. 판별식은
\[
  \Disc(f)=33489=183^2.
\]
따라서 정리 \ref{thm:quartic-galois-decision}에서
\[
  \Gal(f/\Q)\cong A_4.
\]
\end{proof}

\subsection{대칭 사차군: 기약 분해식과 비제곱 판별식}

\begin{example}
\label{ex:quartic-s4}
\[
  f=X^4+X^2+X+1
\]
의 갈루아군은 $S_4$이다.
\end{example}

\begin{proof}
법 $3$으로 줄인
\[
  X^4+X^2+X+1\in\F_3[X]
\]
은 근을 갖지 않고 세 개의 일계수 기약 이차다항식 어느 것으로도 나누어지지 않는다. 따라서 $f$는 $\Q$ 위에서 기약이다.

삼차분해식은
\[
  R_f(Z)=Z^3-2Z^2-3Z+1
\]
이고 가능한 유리근 $\pm1$이 모두 실패하므로 기약이다. 판별식은
\[
  \Disc(f)=257
\]
이며 유리수의 제곱이 아니다. 따라서
\[
  \Gal(f/\Q)\cong S_4.
\]
\end{proof}

\begin{center}
\renewcommand{\arraystretch}{1.22}
\begin{tabular}{llll}
\toprule
$f$ & $R_f$의 형 & $\Disc(f)$ & 갈루아군 \\
\midrule
$X^4-10X^2+1$ & $1+1+1$ & $384^2$ & $V_4$ \\
$X^4-4X^2+2$ & $1+2$ & $2048$ & $C_4$ \\
$X^4-2$ & $1+2$ & $-2048$ & $D_4$ \\
$X^4-7X^2-3X+1$ & $3$ & $183^2$ & $A_4$ \\
$X^4+X^2+X+1$ & $3$ & $257$ & $S_4$ \\
\bottomrule
\end{tabular}
\end{center}

\begin{checkpoint}
\begin{enumerate}[label=(\alph*)]
  \item 예제 \ref{ex:quartic-c4}에서 $\Q(\alpha)$의 비자명한 자동동형 하나를 구체적으로 적어라.
  \item $X^4-2$의 갈루아군 생성원을 근의 순열로 나타내라.
  \item $A_4$ 예제와 $S_4$ 예제에서 삼차분해식의 기약성이 하는 역할을 비교하라.
\end{enumerate}
\end{checkpoint}
